% 原书 p.79–p.81（PDF p.87–p.89）
\section{切触泊松括号（Contact Poisson Bracket）}

\begin{keytakeaways}
\begin{itemize}
  \item 几何图像：切触泊松括号
        $\{F,G\}_{c}=\dd F(X_G)-\dd G(R_\alpha)\cdot F$ 度量 $F$ 沿 $G$ 的
        切触哈密顿流变化有多快，但比辛情形多出修正项 $-\dd G(R_\alpha)F$；
        它\textbf{不是}斜对称的。
  \item 数学工具：经由辛化把切触泊松括号化为辛泊松括号，得到共形协变公式
        (3.37)（证明用 (3.38)–(3.41)）；Legendre 子流形的括号判据
        （习题 3.6.1）。
  \item 本节结论：切触泊松括号在切触同胚下按 (3.37) 共形协变，可用来刻画
        勒让德子流形；其 $C^0$ 下半连续性仍是开问题（问题 3.6.1）。
\end{itemize}
\end{keytakeaways}

% >>> p087
回忆 §2.3 中引入的泊松括号 $\{-,-\}$，用来度量一个输入函数 $F$ 沿另一个
输入函数 $G$ 生成的哈密顿向量场移动时变化有多快。这一量多少有些神秘，却
在意想不到的各种场合出现。在切触几何的框架下，同样可以定义其对应物如下。

% >>> p088
\begin{wdefn}{3.6.1}（参看 [153] 中注 3.5.18；注意符号！）
设 $(M,\xi=\ker\alpha)$ 是切触流形。对任意一对光滑函数
$F,G\in C^\infty(M)$，可定义 $F$ 与 $G$ 的\textbf{切触泊松括号}
（contact Poisson bracket）为
\[
  \{F,G\}_{c}:=\dd F(X_G)-\dd G(R_\alpha)\cdot F,
\]
其中 $X_G$ 是由 $G$ 生成的切触向量场（见注 3.5.3）。
\end{wdefn}

\begin{wrem}{3.6.1}
一般地，切触泊松括号 $\{-,-\}_{c}$ 依赖于 $(M,\xi)$ 的切触 $1$-形式
$\alpha$ 的选取。
\end{wrem}

% [Tutor Note: 易错预警——切触泊松括号与辛泊松括号有一个本质差别：
% $\{F,G\}_{c}$ 一般\textbf{不}斜对称，即
% $\{F,G\}_{c}\neq-\{G,F\}_{c}$（多出的正是
% $\dd G(R_\alpha)F-\dd F(R_\alpha)G$ 之类的项），也不满足 Jacobi 恒等式；
% 它在 $(C^\infty(M),\{\cdot,\cdot\}_{c},\cdot)$ 上给出的是 Leibniz 代数
% 结构，而非 Poisson 代数。这正是原文在定义 3.6.1 处特别提醒「注意符号」
% 的原因。核对 (3.37) 时注意两边的下标 $c$ 与指数因子
% $\ee^{-\kappa(\varphi)}$。]

\begin{wprop}{3.6.1}
设 $\varphi\in\Op{Cont}(M,\xi=\ker\alpha)$ 是共形因子为
$\kappa(\varphi):M\to\R$ 的切触同胚。则对任意
$F,G\in C^\infty(M)$，有下述等式：
\begin{equation}\tag{3.37}
  \bigl\{\ee^{-\kappa(\varphi)}\varphi^*F,\,
  \ee^{-\kappa(\varphi)}\varphi^*G\bigr\}_{c}
  =\ee^{-\kappa(\varphi)}\varphi^*\{F,G\}_{c}.
\end{equation}
\end{wprop}

\begin{wproof}
为证明这一点，我们需要借助辛化（见定义 3.2.1）把切触泊松括号转化为辛几何中
的泊松括号。对给定的切触流形 $(M,\xi=\ker\alpha)$，考虑其辛化
$SM=\R\times M$，辛结构为 $\omega=\dd(\ee^r\alpha)$。对任意切触同胚
$\varphi\in\Op{Cont}(M,\xi)$，可按
\begin{equation}\tag{3.38}
  \widetilde{\varphi}(r,x):=\bigl(r-\kappa(\varphi)(x),\ \varphi(x)\bigr),
  \qquad (r,x)\in\R\times M,
\end{equation}
把它提升为辛同胚 $\widetilde{\varphi}\in\Op{Symp}(SM,\omega)$。对任意函数
$F,G\in C^\infty(M)$，定义
$\widetilde{F}:=\ee^rF$、$\widetilde{G}:=\ee^rG$。由定义 2.3.1，
$\widetilde{F}$ 与 $\widetilde{G}$ 的泊松括号为
\begin{equation}\tag{3.39}
\begin{aligned}
  \{\widetilde{F},\widetilde{G}\}
  &=\dd\widetilde{F}(X_{\widetilde{G}})
   =\ee^r(F\,\dd r+\dd F)\bigl(X_{\ee^rG}\bigr)\\
  &=\ee^r(F\,\dd r+\dd F)\Bigl(X_G-\dd G(R_\alpha)
   \frac{\partial}{\partial r}\Bigr)
   =\ee^r\{F,G\}_{c}.
\end{aligned}
\end{equation}
同时，容易验证（辛）泊松括号在辛同胚作用下不变，这意味着对 (3.38) 中的
$\widetilde{\varphi}$ 有
\begin{equation}\tag{3.40}
  \{\widetilde{\varphi}^*\widetilde{F},
  \widetilde{\varphi}^*\widetilde{G}\}
  =\widetilde{\varphi}^*\{\widetilde{F},\widetilde{G}\}.
\end{equation}
注意
\begin{equation}\tag{3.41}
  (\widetilde{\varphi}^*\widetilde{F})(r,x)
  =\widetilde{F}\bigl(r-\kappa(\varphi)(x),\varphi(x)\bigr)
  =\ee^{r-\kappa(\varphi)(x)}F(\varphi(x))
  =\ee^{-\kappa(\varphi)}(F\circ\varphi)(r,x).
\end{equation}
% >>> p089
于是 (3.39)、(3.40) 与 (3.41) 蕴含
\begin{align*}
  \ee^r\bigl\{\ee^{-\kappa(\varphi)}(F\circ\varphi),\,
  \ee^{-\kappa(\varphi)}(G\circ\varphi)\bigr\}_{c}
  &=\widetilde{\bigl\{\ee^{-\kappa(\varphi)}(F\circ\varphi),\,
   \ee^{-\kappa(\varphi)}(G\circ\varphi)\bigr\}}\\
  &=\{\widetilde{\varphi}^*\widetilde{F},
   \widetilde{\varphi}^*\widetilde{G}\}\\
  &=\widetilde{\varphi}^*\{\widetilde{F},\widetilde{G}\}\\
  &=\widetilde{\varphi}^*\bigl(\ee^r\{F,G\}_{c}\bigr)
   =\ee^{r-\kappa(\varphi)}\{F,G\}_{c}\circ\varphi .
\end{align*}
由此便得到所要的结论 (3.37)。
\end{wproof}

% [Tutor Note: 含糊点拨——(3.39) 中用到的是辛化空间 $SM$ 上
% $\ee^rG$ 的哈密顿向量场公式
% $X_{\ee^rG}=X_G-\dd G(R_\alpha)\partial_r$，它可由
% $\iota_{X_{\ee^rG}}\dd(\ee^r\alpha)=\dd(\ee^rG)$ 分解到
% $\partial_r$、$R_\alpha$、$\xi$ 三个方向逐项验证；这一步教材写作「直接
% 计算」。另外 (3.38) 中 $r$ 的平移量取 $-\kappa(\varphi)$，正是为了让
% $\widetilde\varphi$ 保持 $\omega=\dd(\ee^r\alpha)$，核对时注意符号。]

\begin{wrem}{3.6.2}
事实上，在辛几何中，泊松括号可用来刻画辛流形上的可微映射是否为辛同胚，
办法是验证 $\varphi$ 是否按 (3.40) 的意义保持 $\{-,-\}$。类似地，在切触
几何中，若可微映射 $\varphi$ 保持 $\{-,-\}_{c}$，即对某函数
$\kappa(\varphi)$ 满足 (3.37)，则 $\varphi$ 也是 $(M,\xi=\ker\alpha)$ 上
共形因子为 $\kappa(\varphi)$ 的切触同胚。
\end{wrem}

回忆例 2.3.2 中我们证明了：对辛流形 $(M^{2n},\omega)$ 的（辛）泊松括号，
$n$ 维子流形 $L\subset(M^{2n},\omega)$ 是拉格朗日子流形，当且仅当泊松括号
在函数空间 $I_L=\{F\in C^\infty(M)\mid F|_L=0\}$ 中是封闭的。这一结果有
切触几何的对应物，陈述于下述练习。

\begin{wexer}{3.6.1}（[187] 中命题 1.3）
给定 $(2n+1)$ 维切触流形 $(M,\xi=\ker\alpha)$，设 $L$ 是
$(M,\xi=\ker\alpha)$ 的 $n$ 维子流形。仍记
$I_L:=\{F\in C^\infty(M)\mid F|_L=0\}$。试证明：$L$ 是 $(M,\xi)$ 的
勒让德子流形，当且仅当切触泊松括号 $\{-,-\}_{c}$ 在 $I_L$ 中封闭；换言之，
对任意 $F,G\in I_L$ 有 $\{F,G\}_{c}\in I_L$。
\end{wexer}

在习题 2.3.2 中我们说明了泊松括号不是 $C^0$-连续的，但是
$C^0$-下半连续的（参看定理 2.3.1）。下述切触几何的对应问题至今仍然开
（而与辛情形类似，上半连续性很容易失效）。

\begin{wque}{3.6.1}
给定切触流形 $(M,\xi=\ker\alpha)$，由
$(F,G)\mapsto|\{F,G\}_{c}|_{C^0}$ 定义的泛函
$C^\infty_c(M)\times C^\infty_c(M)\to\R$ 在 $C^0$-拓扑下是否下半连续？
\end{wque}

关于切触泊松括号对切触哈密顿动力学的更多可能应用，见 [81]，特别涉及
勒让德子流形。
